\subsection*{The Complete Dynamic Programming Algorithm}

Now we develop the complete algorithm for the LIS problem.
We follow the recurrence to fill up a DP table, in this case the $L(\cdot)$ array
that stores the length of the optimal solution of all subproblems.
The length of the LIS of the input array $A$ is given by $\max_{1\le k \le n} L(k)$.
The pseudo-code is given below. 
This algorithm runs in $\Theta(n^2)$ time and $\Theta(n)$ space.

\begin{minipage}{0.8\textwidth}
	\aaA {8}{Algorithm LIS-DP~(array $A[1\cdots n]$)}\xxx
	\aaB {5}{for $k = 1 \to n$}\xxx
	\aac {$L[k] = 1$;}\xxx
	\aaC {2}{for $j = 1 \to k - 1$}\xxx
	\aad {if~($A[j] < A[k]$ and $L(j) + 1 > L(k)$): $L[k] = L[j] + 1$;}\xxx
	\aac {end for;}\xxx
	\aab {end for;}\xxx
	\aab {return: $\max_{1\le k \le n} L[k]$;}\xxx
	\aaa {end algorithm;}\xxx
\end{minipage}

Let $k^* = \arg\max_{1\le k \le n} L(k)$, i.e., the LIS of $A$ ends at $A[k^*]$.
We know the length of the LIS of $A$ is $L[k^*]$.
But how to obtain the actual LIS, i.e., $opt(P_{k^*})$? 
First of all, we do not need to explicitly store the actual LIS, i.e., $opt(P_k)$, for all subproblems:
storing each of them explicitly using an array would take $O(n)$ space and
hence storing them all (in this way) takes $O(n^2)$ space, which is prohibitive.
In fact, there is an approach that can store them all in $O(n)$ space!
Why is this possible? Again this is due to the optimal structure property:
the optimal solutions for different subproblems contain duplicate information.

We can use one array $T$ to store \emph{all} optimal solutions, i.e., $\{opt(P_k)\}$. 
To achieve so, 
we define $T[k]$ to store the index/position of the second last element of $opt(P_k)$.
(The last element of $opt(P_k)$ is fixed, i.e., $A[k]$, by the definition of the subproblems.)
See below the $T$ array built for the same example.
%For the same example, this array will be: $$T = [0, 1, 2, 1, 4, 5, 4].$$
\begin{displaymath}
\begin{array}{rllllllllllll}
A & = & [3, &8, &9, &4, &6, &7, &5] \\
opt(P_1) &=& [3]  &&&&&& && T[1] = nul\\
opt(P_2) &=& [3, &8] &&&&&&& T[2] = 1\\
opt(P_3) &=& [3, &8, &9] &&&&&& T[3] = 2\\
opt(P_4) &=& [3, &&&4] &&&&& T[4] = 1\\
opt(P_5) &=& [3, &&&4, &6] &&&& T[5] = 4\\
opt(P_6) &=& [3, &&&4, &6, &7] &&& T[6] = 5\\
opt(P_7) &=& [3, &&&4, &&&5] && T[7] = 4\\
\end{array}
\end{displaymath}

Once $T$ is built~(which we will show),
the optimal solution for \emph{any} $P_k$ can be obtained from $T$ using traceback.
For example, to construct $opt(P_5)$ from $T$, we start with placing $A[5]$
as its last element, by definition. The second last element will be at position/index $T[5] = 4$. 
The element before $A[4]$ must be at position $T[4] = 1$; 
the one right before $A[1]$ is at position $T[1]=nul$ which marks the termination of the traceback.
So, $opt(P_5) = [A[1],A[4],A[5]] = [3, 4, 6]$.
%For example, to construct $opt(P_6)$ from $T$, we start with placing $A[6] = 7$
%as its last element, by definition. The second last element will be at position/index $T[6] = 5$ and the
%corresponding element is $A[5] = 6$. The element before it must be $A[T[5]] = A[4] = 4$;
%the one before is $A[T[4]] = A[1] = 3$. Now $T[1] = nul$ so the traceback terminates,
%which concludes with $opt(P_6) = [3, 4, 6, 7]$.

The pseudo-code for constructing $opt(P_k)$ for \emph{any} $k$ is given below.

\begin{minipage}{0.8\textwidth}
	\aaA {6}{Algorithm LIS-traceback~($A[1\cdots n]$, $T[1\cdots n]$, $k$)}\xxx
	\aaB {3}{while~($k\neq nul$)}\xxx
	\aac {add $A[k]$ to the front of $opt$;}\xxx
	\aac {$k = T[k]$;}\xxx
	\aab {end while;}\xxx
	\aab {return $opt$;}\xxx
	\aaa {end algorithm;}\xxx
\end{minipage}

Of course, to construct the LIS of $A$, i.e., the optimal solution to the original problem,
we can simply call LIS-traceback($A$, $T$, $k^*$), where $k^* = \arg\max_{1\le k \le n} L(k)$.

The only remaining question is to build $T$.
By definition, $T[k]$ stores the index of the second last element of $opt(P_k)$.
Using the same argument for deriving the recurrence for $L(\cdot)$, 
the second last element of $opt(P_k)$
is $A[j]$, where $1\le j < k$ and $A[j] < A[k]$, that maximizes $L(j)$.
Formally, $$\textstyle T[k] = \arg\max_{1\le j < k, A[j] < A[k]} L(j).$$

Note the connection between $T$ and $L$.
Not a surprise, $T$ can be constructed on the fly along with constructing $L$.
Whenever $L[k]$ gets updated, meaning a larger $L[j]$, where $1\le j \le k$ and $A[j] < A[j]$, is found, 
we update $T[k]$ accordingly to keep track of $j$, i.e., the index of currently the best second last element.
The pseudo-code for constructing $T$ (and $L$) is given below.

\begin{minipage}{0.8\textwidth}
	\aaA {9}{Algorithm LIS-DP-T~(array $A[1\cdots n]$)}\xxx
	\aaB {6}{for $k = 1 \to n$}\xxx
	\aac {$L[k] = 1$;}\xxx
	\aac {$T[k] = nul$;}\xxx
	\aaC {2}{for $j = 1 \to k - 1$}\xxx
	\aad {if~($A[j] < A[k]$ and $L(j) + 1 > L(k)$): $L[k] = L[j] + 1$ and $T[k] = j$}\xxx
	\aac {end for;}\xxx
	\aab {end for;}\xxx
	\aab {return $T$;}\xxx
	\aaa {end algorithm;}\xxx
\end{minipage}


%%The pseudo-code is given above.
%%We maintain two arrays, $L(k)$, to store the length of the $opt(k)$,
%%and $T(k)$, to maintain the tracking-back pointers, $1 \le k \le n$.
%%The optimal solution $opt(k)$ can be constructed by using $T$.
%%The running time of the algorithm is $\Theta(n^2)$.


\subsection*{An $\Theta(n\log n)$ Algorithm}

The LIS problem can be solved in $\Theta(n\log n)$ time.
The idea is to maintain the ``best'' IS of each length, so that
longer IS can be built on top of them when the next number gets processed.
Here the ``best'' means the subsequence ends with the smallest possible number.
This is intuitive as it facilitates further extension.

Formally, we define $S_k[l]$ as the smallest number in $A[1\cdots k]$ such
that there exists an IS of length $l$ that ends at this number.
It helps to work on an example to fully understand above defintion.

$A = [3, 8, 5, 9, 6, 2, 7, 4]$. Then we have:\\
$S_1 = [3]$; \\
$S_2 = [3, 8]$; \\
$S_3 = [3, 5]$; \\
$S_4 = [3, 5, 9]$; \\
$S_5 = [3, 5, 6]$; \\
$S_6 = [2, 5, 6]$; \\
$S_7 = [2, 5, 6, 7]$; \\
$S_8 = [2, 4, 6, 7]$.

It is important to keep in mind that $S_k[l]$ corresponds to an IS of length $l$.
$S_k[l]$ does not store the entire IS but just its last element.

We observe that, for every $k$, the array $S_k$ is sorted in ascending order.
This is not an coinsidence. We now prove it, by contradiction.
Suppose conversely that, for some $k$, $S_k[i] > S_k[j]$ where $i < j$.
By the definition of $S_k[j]$, there exists an IS
of length $j$ that ends at $S_k[j]$. Let $a_i$ be the $i$-th number
in this IS; we have $a_i < S_k[j]$.
Note that $a_i$ is also the last number of an IS of length $i$.
Since, by definition, $S_k[i]$ stores the \emph{smallest}
last number across all length-$i$ ISs,
we know that $S_k[i] \le a_i$. This results in a contradiction.

Having $S_k$ is increasing is critical for obtaining the faster algorithm,
as it allows for binary search which takes $O(\log n)$ time. 

Suppose that we are able to build $S_k$, for all $k = 1, 2, \cdots, n$,
what is the LIS of $A$? Its length will be $|S_n|$; 
the actual LIS can be constructed using traceback as we will see shortly.

The remaining question is to build $\{S_k\}$. We do it gradually,
i.e., to build $S_{k+1}$ based on $S_k$.
You might also have observed in above example that, at most one element
needs to be updated. This is indeed true.
The updating procedure is quite simple: if $A[k+1]$ is larger than the last element of $S_k$~(or, $S_k$ is empty),
we append $A[k+1]$ to its end; otherwise, we locate~(using binary search) the smallest $i$
such that $S_k[i] > A[k+1]$ and then replace $S_k[i]$ with $A[k+1]$.
Please take a moment to verify that this procedure works by applying to above example.

Now we formally prove that this procedure results in the correct $S_{k+1}$.
In the case that $A[k+1]$ is larger than the last element of $S_k$, or $S_k$ is empty,
we know that a longer IS, of length $|S_k| + 1$, can be formed, by concatenating
the IS implied by the last element of $S_k$ and $A[k+1]$.
%the IS of $|S_k|$~(whose last element is the last element of $S_k$) and $A[k+1]$.
Also, $A[k+1]$ must be used to form an IS of length $|S_k| + 1$ because
without it the length of the LIS in $A[1\cdots k]$ is $|S_k|$. Hence, $A[k+1]$ is
also the smallest possible last element, because it is the only choice.
We also have $S_{k+1}[l] = S_k[l]$ for each $l = 1, 2, \cdots, |S_k|$, since
the best length-$l$ IS in $A[1\cdots k+1]$ will not use $A[k+1]$ because $S_k[l] < A[k+1]$.
These analysis combined proves that simply appending $A[k+1]$ to $S_k$ obtains $S_{k+1}$.

Now consider the case that $S_k[i] > A[k+1]$, where $i$ is the smallest such index.
This implies that $S_k[l] < A[k+1]$ for all $1 \le l < i$.
Using above analysis, we have that $S_{k+1}[l] = S_k[l]$, for $1 \le l < i$.
%then the first $i-1$ elements in $S_k$ should remain unchanged as they are all smaller than $A[k+1]$.
Next, $S_{k+1}[i]$ must be $A[k+1]$,
%$S_k[i]$ should be replaced with $A[k+1]$ 
as now we have a ``better'' IS of length $i$,
by concatenating the IS of length $(i-1)$ implicitly stored in $S_k[i-1]$ and $A[k+1]$; 
this new length-$i$ IS is better because $A[k+1] < S_k[i]$.
Last, $S_{k+1}[l] = S_k[l]$ for $l = i + 1, \cdots, |S_k|$.
This can be proved by contradiction.
Assume that $S_{k+1}[l] = A[k+1]$ for some $l \ge i + 1$. 
Consider the IS of length $l$ it implies.
The $i$-th element in it, which we call $a_i$, is the last element of an IS of length $i$,
so we must have $S_k[i] \le a_i < A[k+1]$ which is a contradiction.
All these combined proves that, for this case, simply replacing $S_k[i]$ with $A[k+1]$ results in the correct $S_{k+1}$.

The pseudo-code is given below.  Note that we will not maintain separate arrays $S_k$; 
we only maintain one array $S$ and update it as $k$ increases.

\begin{minipage}{0.8\textwidth}
	\aaA {7}{Algorithm faster-LIS~(array $A[1\cdots n]$)}\xxx
	\aab {$S=\emptyset$;}\xxx
	\aaB {3}{for $k = 1 \to n$}\xxx
	\aac {if $|S| = 0$ or $A[k] > S[|S|]$:  append $A[k]$ to $S$;}\xxx
	\aac {else: binary-search($S$, $A[k]$) to locate the smallest $i$ such that $S[i] > A[k]$, then $S[i] = A[k]$;}\xxx
	\aab {end for;}\xxx
	\aab {return $|S|$;}\xxx
	\aaa {end algorithm;}\xxx
\end{minipage}

Above algorithm computes the length of the LIS. To construct the actual LIS, we need two modifications.
First, instead of storing the numbers in $A$, we use $S$ to store the corresponding indexes.
Second, we maintain tracing back pointers stored in an array $T$.

\begin{minipage}{0.8\textwidth}
	\aaA {18}{Algorithm faster-LIS-full~(array $A[1\cdots n]$)}\xxx
	\aab {$S=\emptyset$; }\xxx
	\aaB {9}{for $k = 1 \to n$}\xxx
	\aaC {3}{if $|S| = 0$ or $A[k] > S[|S|]$}\xxx
	\aad {$T[k] = S[|S|]$ or $T[k] = null$ if $|S|=0$;}\xxx
	\aad {append $k$ to $S$;}\xxx
	\aaC {4}{else}\xxx
	\aad {binary-search to locate the smallest $i$ such that $A[S[i]] > A[k]$;}\xxx
	\aad {$S[i] = k$;}\xxx
	\aad {$T[k] = S[i-1]$ or $T[k] = nul$ if $i=1$;}\xxx
	\aac {end if;}\xxx
	\aab {end for;}\xxx
	\aab {$j = S[|S|]$; $opt=\emptyset$;}\xxx
	\aaB {3}{while~($j\neq nul$)}\xxx
	\aac {add $A[j]$ to the front of $opt$;}\xxx
	\aac {$j = T[j]$;}\xxx
	\aab {end while;}\xxx
	\aab {return $opt$;}\xxx
	\aaa {end algorithm;}\xxx
\end{minipage}

Both algorithms runs in $\Theta(n\log n)$ time as binary-search takes $\Theta(\log n)$ time.
